\subsection{Sparse High-Dimensional Inference}

High-dimensional statistics studies regimes where the number of candidate parameters is comparable to, or much larger than, the sample size.  Estimation is possible only after adding structure.

\begin{definition} \label{def:sparse-vector-and-high-dimensional-regime} \hlt{Sparsity and High Dimension}\\
A vector $\theta\in\R^p$ is \hlt{$s$-sparse} if at most $s$ of its coordinates are nonzero.  A sequence of problems is \hlt{high-dimensional} when $p=p_n$ grows with $n$, often with $p_n\gg n$.
\end{definition}

\begin{definition} \label{def:lasso-estimator} \hlt{Lasso}\\
For linear regression $y=X\beta+\varepsilon$, the \hlt{lasso} estimator solves
\[
\widehat{\beta}_\lambda
\in\operatorname*{arg\,min}_{\beta\in\R^p}
\left\{\frac{1}{2n}\norm{y-X\beta}_2^2+\lambda\norm{\beta}_1\right\}.
\]
The $\ell^1$ penalty encourages sparse solutions and is convex.
\end{definition}

\begin{definition} \label{def:restricted-eigenvalue-condition} \hlt{Restricted Eigenvalue Condition}\\
Let $S\subseteq\{1,\ldots,p\}$ with $|S|=s$.  A design matrix satisfies a \hlt{restricted eigenvalue condition} on $S$ if there exists $\kappa>0$ such that
\[
\frac{\norm{Xv}_2}{\sqrt{n}}\geq \kappa\norm{v_S}_2
\]
for all $v$ in a cone such as $\norm{v_{S^c}}_1\leq3\norm{v_S}_1$.
\end{definition}

\begin{theorem} \label{thm:lasso-basic-rate-form} \hlt{Basic Lasso Rate Form}\\
Assume the true coefficient vector is $s$-sparse, the design satisfies a restricted eigenvalue condition, and $\lambda$ dominates $\norm{X^\top\varepsilon/n}_\infty$.  Then the lasso obeys bounds of the form
\[
\frac{1}{n}\norm{X(\widehat{\beta}_\lambda-\beta_0)}_2^2
\lesssim s\lambda^2,
\qquad
\norm{\widehat{\beta}_\lambda-\beta_0}_1\lesssim s\lambda.
\]
\end{theorem}

\begin{definition} \label{def:false-discovery-rate} \hlt{False Discovery Rate}\\
For multiple testing, let $V$ be the number of false rejections and $R$ the total number of rejections.  The \hlt{false discovery rate} is
\[
\operatorname{FDR}=\mathbb{E}\!\left[\frac{V}{\max(R,1)}\right].
\]
It controls the expected false-rejection proportion among discoveries.
\end{definition}

\begin{remark} \label{rem:high-dimensional-cross-reference-optimization} \hlt{Statistical and Computational Structure}\\
High-dimensional procedures combine probability concentration, convex optimization, and model-selection error control.  Their computational side belongs naturally with the optimization and numerical methods of Part 14.
\end{remark}
